\section{Replacement-covariance factorization boundary}
\textbf{Status: PROVED/EXACT within the declared transformation model; P versus NP remains OPEN.}
Let \(\mathcal P\) be any nonempty class of exact program codes, let
\(\mathcal W\) be program transformations, and let
\(X:\mathcal P\times\mathcal F\to\mathcal O\) satisfy
\[
 X(W(A),F)=X(A,F)
 \quad(W\in\mathcal W,\ A\in\mathcal P,\ F\in\mathcal F).
\]
Then \(X(\cdot,F)\) is constant on every component of the transformation
graph. If \(\mathcal W\) contains each constant replacement
\(\kappa_B(A)=B\), then
\[
 X(A,F)=R(F)
\]
for a source-only functional \(R\). Thus a universal resource cannot be both
implementation-sensitive and invariant under every exact polynomial
language-preserving transformation. Excluding replacements preserves
intensional information only at the price of a new coverage theorem for
arbitrary SAT implementations. Source factorization does not imply that the
original execution materializes or pays for \(R(F)\).

Applied to a compiler landing in a fixed FLMV-forbidden ordinary class
\(\mathcal G\), full replacement covariance reduces the desired statement to
\(\SAT\in P\Rightarrow\SAT\in\mathcal G\), already a
separation-strength endpoint rather than an intermediate obligation. The
finite audit exhaustively checked 586,010 labelings and 1,177 assertions,
including Horn, 2SAT, affine/Tseitin, bounded-treewidth, independent-product,
padding, encoding and DIRECT/DEAD/COPY/DELAY/INTERP/RECOMP/REPLACE controls.
The next OPEN target is \texttt{TYPED-COVARIANCE-COVERAGE1}.
See \path{research/review/reports/model_bridge/2026-10-11_1800_MODELBRIDGE-0016.tex}.

